\scanpage{202}
\mspage{11.}
Soit $(\mathscr R,r)$ une donnée de ramification sur $S$. Une famille d'objets $X_i/S_i$ de $\mathscr R$ ($S_i\in\mathrm{Ob}\,Et(S)$) est dite \uline{génératrice} si pour tt objet $X$ de $\mathscr R$ sur un $S'$ de $Et(S)$, les $S''$ étales sur $S$ tels qu'il existe $i$ et $S''\unsure{-|\to}S_i$ avec $X_{S''}$ subordonné à $X_{i\,S''}$ \inserted{et $X_{i\,S''}\deleted{\unsure{\to}}\to e_{S''}$ \deleted{\unsure{étant}} couvrant}, recouvrent $S'$.
\diagram{\begin{tikzcd}[ampersand replacement=\&]
S\arrow[dr,dash]\&{}\&S_i\\
{}\&S'\arrow[dr,dash]\&{}\\
{}\&{}\&S''
\end{tikzcd}}
Si \inserted{de plus} les objets $X_i$ sont galoisiens dans $\mathscr R_{S_i}$, \deleted{\unsure{sur les}} de groupe $G_i$, d'où $(Z_i,G_i)$ objet à opérateurs $r(\deleted{\unsure{P|R}}X_i)$ sur $S_i$, les $Z_i,G_i$ doivent permettre de \uline{reconstruire} $(\mathscr R,r)$, une fois précisé des données supplémentaires permettant de recoller ; mais la formulation de telles données est pénible et finalement peu utile.

Il est plus commode de supposer donnée déjà une \uline{catégorie fibrée} $C$ sur \inserted{$Et$}$(S)$, et un foncteur \inserted{cartésien} $\varphi:C\longrightarrow\mathscr R$. On dit que $\varphi$ est générateur si les $\varphi(M)$ ($M\in\mathrm{Ob}\,C$) forment une famille génératrice, i.e si pour tt $X\in\mathrm{Ob}\,\mathscr R$ sur un $S'\in\mathrm{Ob}\,Et(S)$, les $S''$ étales sur $S'$ tels que $X_{S'}$ soit subordonné à un objet \inserted{$M$} de $C_{S''}$ \inserted{avec $\varphi(M)\to e_{S^{\unsure{\prime|\prime\prime}}}$ couvrant,} forment un recouvrement de $S'$. Posant $s=r\varphi$
\[
s:C\unsure{-|\longrightarrow}\mathrm{Fl}_S
\]
je dis qu'on peut alors, à équivalence près,
% archive sheet number: 201
